\documentclass[reqno,11pt]{amsart}
\usepackage[left=1in,right=1in,top=1in,bottom=1in]{geometry}
\setcounter{tocdepth}{1}
\usepackage{amsmath}
\usepackage{amssymb}
\usepackage{mathtools}
\usepackage{graphicx}
\usepackage{tikz}
\usepackage{esint}
\usepackage{xcolor}
\usetikzlibrary{arrows, decorations.markings,matrix,trees}
\usepackage{hyperref}
\hypersetup{
   colorlinks=true,
   citecolor=blue,
   filecolor=blue,
   linkcolor=blue,
   urlcolor=blue
}

\newcommand{\al}{\alpha}
\newcommand{\vt}{\vartheta}
\newcommand{\lda}{\lambda}
\newcommand{\om}{\Omega}
\newcommand{\pa}{\partial}
\newcommand{\va}{\varepsilon}
\newcommand{\ud}{\mathrm{d}}
\newcommand{\be}{\begin{equation}}
\newcommand{\ee}{\end{equation}}
\newcommand{\w}{\omega}
\newcommand{\Lda}{\Lambda}

\newcommand{\A}{\mathbf{A}}
\newcommand{\cA}{\mathcal{A}}
\newcommand{\bA}{\mathbb{A}}

\newcommand{\B}{\mathbb{B}}
\newcommand{\bb}{\mathbf{b}}
\newcommand{\cB}{\mathcal{B}}

\newcommand{\CC}{\mathbf{C}}
\newcommand{\cC}{\mathcal{C}}
\providecommand{\C}{}
\renewcommand{\C}{\mathbb{C}}

\newcommand{\bd}{\mathbf{d}}
\newcommand{\bD}{\mathbb{D}}
\newcommand{\rD}{\mathrm{D}}
\newcommand{\cD}{\mathcal{D}}

\newcommand{\E}{\mathbb{E}}
\newcommand{\cE}{\mathcal{E}}
\newcommand{\e}{\mathbf{e}}
\newcommand{\re}{\mathrm{e}}

\newcommand{\fF}{\mathbf{F}}
\newcommand{\cF}{\mathcal{F}}

\newcommand{\bG}{\mathbb{G}}
\newcommand{\fg}{\mathbf{G}}
\newcommand{\cG}{\mathcal{G}}

\newcommand{\rh}{\mathrm{h}}
\newcommand{\rH}{\mathrm{H}}
\newcommand{\bH}{\mathbb{H}}

\newcommand{\I}{\mathbf{I}}
\newcommand{\ii}{\mathrm{i}}

\newcommand{\jj}{\mathrm{j}}

\newcommand{\kk}{\mathrm{k}}

\newcommand{\LL}{\mathbf{L}}
\newcommand{\cL}{\mathcal{L}}

\newcommand{\Z}{\mathbb{Z}}

\newcommand{\M}{\mathcal{M}}
\newcommand{\MM}{\mathbb{M}}
\newcommand{\bM}{\mathbf{M}}
\newcommand{\m}{\mathbf{m}}

\newcommand{\cN}{\mathcal{N}}
\newcommand{\N}{\mathbb{N}}
\newcommand{\rN}{\mathrm{N}}
\newcommand{\n}{\mathbf{n}}

\newcommand{\OO}{\mathbb{O}}

\newcommand{\cP}{\mathcal{P}}
\newcommand{\PP}{\mathbf{P}}
\newcommand{\p}{\mathbf{p}}

\newcommand{\Q}{\mathbf{Q}}

\newcommand{\R}{\mathbb{R}}
\newcommand{\RR}{\mathbf{R}}
\newcommand{\cR}{\mathcal{R}}

\newcommand{\cS}{\mathcal{S}}
\newcommand{\Ss}{\mathbb{S}}
\newcommand{\fS}{\mathbf{S}}

\newcommand{\cT}{\mathcal{T}}

\newcommand{\uu}{\mathbf{u}}
\newcommand{\U}{\mathbf{U}}

\newcommand{\V}{\mathbf{V}}
\newcommand{\fv}{\mathbf{v}}
\newcommand{\cV}{\mathcal{V}}

\newcommand{\W}{\mathbf{W}}
\newcommand{\cW}{\mathcal{W}}

\newcommand{\X}{\mathbf{X}}
\newcommand{\cX}{\mathcal{X}}

\newcommand{\Y}{\mathbf{Y}}
\newcommand{\cY}{\mathcal{Y}}

\newcommand{\wc}{\rightharpoonup}
\newcommand{\Sn}{\Ss^n}
\newcommand{\HH}{\mathcal{H}}
\newcommand{\FF}{\mathcal{F}}

\newcommand{\vp}{\varphi}
\newcommand{\da}{\downarrow}
\newcommand{\T}{\mathrm{T}}
\newcommand{\ga}{\gamma}
\newcommand{\Ga}{\Gamma}
\newcommand{\ep}{\varepsilon}
\newcommand{\sg}{\sigma}
\newcommand{\ift}{\infty}
\newcommand{\wt}{\widetilde}
\newcommand{\wh}{\widehat}
\newcommand{\f}{\frac}
\newcommand{\D}{\mathrm{D}}
\newcommand{\ol}{\overline}
\newcommand{\Ra}{\Rightarrow}
\newcommand{\op}{\operatorname}
\newcommand{\Sg}{\Sigma}
\newcommand{\nn}{\nonumber}
\newcommand{\na}{\nabla}
\newcommand{\dint}{\displaystyle\int}
\newcommand{\dlim}{\displaystyle\lim}
\newcommand{\dsup}{\displaystyle\sup}
\newcommand{\dmin}{\displaystyle\min}
\newcommand{\dmax}{\displaystyle\max}
\newcommand{\dinf}{\displaystyle\inf}
\newcommand{\dsum}{\displaystyle\sum}
\newcommand{\abs}[1]{\lvert#1\rvert}
\newcommand{\sbt}{\,\begin{picture}(-1,1)(-1,-3)0le*{3}\end{picture}\ }

%\renewcommand{\labelenumi}{(\roman{enumi})}

\DeclareMathOperator{\dist}{dist}
\DeclareMathOperator{\avg}{avg}
\DeclareMathOperator{\diam}{diam}
\DeclareMathOperator{\BMO}{BMO}
\DeclareMathOperator{\VMO}{VMO}
\DeclareMathOperator{\sgn}{sgn}
\DeclareMathOperator{\supp}{supp}
\DeclareMathOperator{\im}{im}
\DeclareMathOperator{\Lip}{Lip}
\DeclareMathOperator{\Vol}{Vol}
\DeclareMathOperator{\tr}{tr}
\DeclareMathOperator*{\osc}{osc}
\DeclareMathOperator{\Length}{Length}
\DeclareMathOperator{\diag}{diag}
\DeclareMathOperator{\Dist}{Dist}
\DeclareMathOperator{\graph}{graph}
\DeclareMathOperator{\even}{even}
\DeclareMathOperator{\odd}{odd}
\DeclareMathOperator{\sing}{sing}
\DeclareMathOperator{\reg}{reg}
\DeclareMathOperator{\rank}{rank}
\DeclareMathOperator{\const}{const}
\DeclareMathOperator{\loc}{loc}
\DeclareMathOperator{\argmin}{argmin}
\DeclareMathOperator{\ad}{ad}
\DeclareMathOperator{\id}{id}
\DeclareMathOperator{\Int}{Int}
\DeclareMathOperator{\fun}{fun}
\DeclareMathOperator{\curl}{curl}
\DeclareMathOperator{\vol}{vol}
\DeclareMathOperator*{\esssup}{ess\,sup}
\DeclareMathOperator*{\essinf}{ess\,inf}
\DeclareMathOperator*{\essrange}{ess\,range}

\providecommand{\fint}{\mathop{\ooalign{$\int$\cr\hidewidth$-$\hidewidth\cr}}\nolimits}
\newenvironment{addedblock}{\begingroup\color{blue}}{\endgroup}

\def\<{\langle}\def\>{\rangle}
\def\({\left(}\def\){\right)}
\numberwithin{equation}{section}
\theoremstyle{plain}
\newtheorem{thm}{Theorem}[section]
\newtheorem{cor}[thm]{Corollary}
\newtheorem{con}[thm]{Conjecture}
\newtheorem{lem}[thm]{Lemma}
\newtheorem{prop}[thm]{Proposition}
\newtheorem{assum}[thm]{Assumption}

\theoremstyle{definition}
\newtheorem{defn}[thm]{Definition}
\newtheorem{example}[thm]{Example}
\newtheorem{q}[thm]{Question}

\theoremstyle{remark}
\newtheorem{rem}[thm]{Remark}

\title[Lean statements for stationary Sobolev monotonicity]{Lean Statements for the Stationary Sobolev Monotonicity Formalization}

\begin{document}

\maketitle
\tableofcontents

\section{Conventions}

This document is a theorem-and-definition ledger for the Lean development.
Proofs are intentionally omitted.  Each theorem, proposition, lemma, or
definition heading records the source Lean file and the Lean declaration name.

Lean's \texttt{Domain n} is read as $\R^n$, \texttt{Target m} as $\R^m$,
and \texttt{Gradient n m} as the space of arrays $G=(G_i)_{i=1}^n$ with
$G_i\in\R^m$.  If $G:\R^n\to(\R^m)^n$, then
\[
  \Theta(G,a,\rho)
  :=
  \rho^{2-n}\int_{B(a,\rho)} \abs{G(x)}^2\,\ud x
\]
denotes the mathematical version of Lean's \texttt{weakTheta G a rho}.
The annular radial-energy term
\[
  \mathcal R(G,a;s,r)
  :=
  2\int_{B(a,r)\setminus B(a,s)}
    \abs{x-a}^{2-n}\abs{\partial_{r,a}G(x)}^2\,\ud x
\]
denotes the mathematical version of Lean's
\texttt{weakMonotonicityRhs G a s r}.

\section{Main Public Theorems}

\begin{thm}[\texttt{MainTheorem.lean}: \texttt{stationaryW12LocMonotonicityFormula\_euclidean}]
\label{thm:stationaryW12LocMonotonicityFormulaEuclidean}
Let $n,m\in\N$ with $n\neq 0$.  Let $\om\subset\R^n$ be a measurable set,
let $u:\R^n\to\R^m$, and let
\[
  \mathcal U : \texttt{StationaryW12LocMap}\; u\; \om
\]
be a stationary local $W^{1,2}$ witness with displayed weak gradient
\[
  G_{\mathcal U}:=\mathcal U.\texttt{w12.weakGrad}.
\]
Assume that $a\in\R^n$, $R_0,s,r\in\R$ satisfy
\[
  \ol{B(a,R_0)}\subset\om,\qquad 0<s,\qquad s<r,\qquad r<R_0.
\]
Then the normalized weak-energy increment is exactly the annular radial-energy
term:
\[
  \Theta(G_{\mathcal U},a,r)-\Theta(G_{\mathcal U},a,s)
  =
  \mathcal R(G_{\mathcal U},a;s,r).
\]
\end{thm}

\begin{thm}[\texttt{MainTheorem.lean}: \texttt{stationaryW12LocMonotonicity\_euclidean}]
\label{thm:stationaryW12LocMonotonicityEuclidean}
Let $n,m\in\N$ with $n\neq 0$.  Let $\om\subset\R^n$ be a measurable set,
let $u:\R^n\to\R^m$, and let
\[
  \mathcal U : \texttt{StationaryW12LocMap}\; u\; \om
\]
be a stationary local $W^{1,2}$ witness with displayed weak gradient
\[
  G_{\mathcal U}:=\mathcal U.\texttt{w12.weakGrad}.
\]
Assume that $a\in\R^n$, $R_0,s,r\in\R$ satisfy
\[
  \ol{B(a,R_0)}\subset\om,\qquad 0<s,\qquad s\le r,\qquad r<R_0.
\]
Then the normalized weak energy is monotone in the radius:
\[
  \Theta(G_{\mathcal U},a,s)
  \le
  \Theta(G_{\mathcal U},a,r).
\]
\end{thm}

\begin{rem}[\texttt{MainTheorem.lean}: theorem inputs]
The Lean package \texttt{StationaryW12LocMap u Omega} contains the local
$L^2$ data for $u$, the displayed weak gradient, local $L^2$ and a.e.
measurability data for that gradient, the distributional weak-gradient
identity, and the vanishing domain first-variation identity.  The theorem does
not assume any target-manifold constraint.
\end{rem}

\section{Planned File-by-File Transcription}

The remaining entries will be inserted one Lean file at a time.  Within each
file section, every public definition and every proved declaration
(\texttt{def}, \texttt{abbrev}, \texttt{structure}, \texttt{theorem},
\texttt{lemma}, \texttt{instance}, and related named conclusions) will receive
a corresponding LaTeX entry, with no proof text.

\subsection{Public API layer}

\subsubsection*{\texttt{StationaryHarmonicMap/Basic.lean}}
\subsubsection*{\texttt{StationaryHarmonicMap/Euclidean.lean}}
\subsubsection*{\texttt{StationaryHarmonicMap/WeakStationarity.lean}}
\subsubsection*{\texttt{StationaryHarmonicMap/WeakGradientBridge.lean}}
\subsubsection*{\texttt{StationaryHarmonicMap/L2LocBridge.lean}}
\subsubsection*{\texttt{StationaryHarmonicMap/SobolevBridge.lean}}
\subsubsection*{\texttt{StationaryHarmonicMap/FirstVariationBridge.lean}}
\subsubsection*{\texttt{StationaryHarmonicMap/StationarityBridge.lean}}
\subsubsection*{\texttt{StationaryHarmonicMap/SobolevWitness.lean}}
\subsubsection*{\texttt{StationaryHarmonicMap/StationaryMap.lean}}
\subsubsection*{\texttt{StationaryHarmonicMap/MainTheorem.lean}}
\subsubsection*{\texttt{StationaryHarmonicMap/API.lean}}
\subsubsection*{\texttt{Examples/UseMainTheorem.lean}}
\subsubsection*{\texttt{Examples/StationaryMonotonicity.lean}}

\subsection{Radial geometry and cutoff layer}

\subsubsection*{\texttt{StationaryHarmonicMap/EnergyQuantities.lean}}
\subsubsection*{\texttt{StationaryHarmonicMap/RadialGeometry.lean}}
\subsubsection*{\texttt{StationaryHarmonicMap/RadialCutoffs.lean}}
\subsubsection*{\texttt{StationaryHarmonicMap/PrimitiveCutoffs.lean}}
\subsubsection*{\texttt{StationaryHarmonicMap/RadialIntegrability.lean}}
\subsubsection*{\texttt{StationaryHarmonicMap/RadialIdentity.lean}}

\subsection{Radial measure and radius calculus layer}

\subsubsection*{\texttt{StationaryHarmonicMap/RadialMeasure.lean}}
\subsubsection*{\texttt{StationaryHarmonicMap/RadiusWeights.lean}}
\subsubsection*{\texttt{StationaryHarmonicMap/RadiusPrimitive.lean}}
\subsubsection*{\texttt{StationaryHarmonicMap/BallIntegralAC.lean}}
\subsubsection*{\texttt{StationaryHarmonicMap/RadiusWeightedDerivative.lean}}
\subsubsection*{\texttt{StationaryHarmonicMap/RadiusWeightedRepresentation.lean}}
\subsubsection*{\texttt{StationaryHarmonicMap/RadiusAnalysis.lean}}
\subsubsection*{\texttt{StationaryHarmonicMap/RadiusFormulas.lean}}

\subsection{Boundary identity and monotonicity layer}

\subsubsection*{\texttt{StationaryHarmonicMap/BoundaryBasics.lean}}
\subsubsection*{\texttt{StationaryHarmonicMap/BoundaryFromRadial.lean}}
\subsubsection*{\texttt{StationaryHarmonicMap/MonotonicityFinal.lean}}
\subsubsection*{\texttt{StationaryHarmonicMap/MonotonicityRoutes.lean}}
\subsubsection*{\texttt{StationaryHarmonicMap/MonotonicityEuclidean.lean}}
\subsubsection*{\texttt{StationaryHarmonicMap/CenterTranslation.lean}}
\subsubsection*{\texttt{StationaryHarmonicMap/Monotonicity.lean}}

\end{document}
